% ============================================================
% AUTOMATISCH GEGENEREERD BESTAND
% Niet handmatig aanpassen.
% ============================================================

\ifdefined\HCode
  % Geen afzonderlijke module-openingspagina in HTML.
\else
  \FVDmoduleopening
    {Van krachten tot kernen}
    {In deze cursus bouwen we verder op de fysica uit de tweede graad en onderzoeken we hoe krachten en velden ons helpen om verschijnselen op heel verschillende schalen te begrijpen: van de beweging van alledaagse voorwerpen en planeten tot de wisselwerking tussen elektrische ladingen en de processen in de atoomkern. We starten met de gereedschapskist van de fysicus. Grootheden en eenheden, beduidende cijfers, wetenschappelijke notatie en vectoren vormen de taal waarmee we fysische verschijnselen nauwkeurig kunnen beschrijven. Daarna hernemen en verdiepen we de wetten van Newton. We leren krachten analyseren, vectorieel voorstellen en gebruiken om de beweging van voorwerpen te verklaren en te voorspellen. Vervolgens maken we de stap van krachten die we rechtstreeks waarnemen naar krachten die op afstand werken. Vanuit de gravitatie onderzoeken we hoe massa's elkaar aantrekken en hoe het begrip veld toelaat om die wisselwerking te beschrijven. Daarna ontdekken we dat elektrische ladingen op een verrassend gelijkaardige manier met elkaar interageren. De wet van Coulomb, het elektrische veld, elektrische potentiële energie, potentiaal en spanning bouwen samen een steeds rijker model van elektrische wisselwerking op. Vanuit elektrische ladingen zetten we de stap naar magnetisme. We onderzoeken hoe bewegende ladingen en elektrische stromen magnetische velden veroorzaken en hoe magnetische velden op bewegende ladingen kunnen inwerken. Daarna ontdekken we met elektromagnetische inductie dat een veranderend magnetisch veld op zijn beurt elektrische effecten kan veroorzaken. Zo groeien elektriciteit en magnetisme uit van twee ogenschijnlijk verschillende verschijnselen tot onderdelen van één samenhangend fysisch verhaal. Ten slotte richten we onze blik op een veel kleinere schaal: de atoomkern. We onderzoeken de opbouw en stabiliteit van kernen, radioactief verval en kernreacties en bestuderen hoe massa en energie met elkaar verbonden zijn. Begrippen zoals halveringstijd, bindingsenergie en massadefect laten zien hoe dezelfde fysische manier van denken ook processen kan beschrijven die zich diep binnenin de materie afspelen. Doorheen de cursus gebruiken we wiskunde als taal van de fysica. Vectoren helpen ons krachten en velden te beschrijven, functies tonen hoe fysische grootheden van elkaar afhangen en exponentiële modellen krijgen een concrete betekenis bij radioactief verval. We leren niet alleen formules gebruiken, maar ook modellen opbouwen, verbanden herkennen en resultaten kritisch interpreteren. Experimenten, metingen en gegevensanalyse vormen daarbij een essentieel onderdeel van de fysica. We leren onderzoeksvragen stellen, metingen uitvoeren, gegevens verwerken en nagaan in welke mate een model de werkelijkheid beschrijft. Ook toepassingen uit wetenschap en technologie laten zien hoe de onderzochte principes buiten het klaslokaal worden gebruikt. Fysica wordt zo geen verzameling afzonderlijke formules en wetten. Van Newtons bewegingswetten over gravitatie en elektromagnetisme tot processen in de atoomkern zoeken we telkens naar dezelfde onderliggende vragen: welke wisselwerkingen spelen een rol, hoe kunnen we ze modelleren en wat kunnen we ermee verklaren en voorspellen?}
    {%
    \FVDmodulethemetile{1}{Gereedschapskist en Newton - Van meten en vectoren tot krachten en beweging}
    \FVDmodulethemetile{2}{Krachten op afstand - Van gravitatie tot elektrische velden}
    \FVDmodulethemetile{3}{Elektriciteit wordt magnetisme - Van bewegende ladingen tot elektromagnetische inductie}
    \FVDmodulethemetile{4}{Binnen in de kern- Van kernstructuur tot radioactief verval}
    }
\fi
